% ------------------------------------------------------------
% Page 100
% ------------------------------------------------------------

\textbf{Le Maire : M. Fages}

Peu après 14 heures, le Maire, pharmacien de son état, me reçoit dans son officine ; je
me présente. Serrement de mains, mais pas un mot de bienvenue, que j’espérais
quand même de sa part. Je lui donne mon arrêté de nomination ; il le prend, se retire
dans son bureau (l’arrière boutique), revient presque aussitôt, et me rend le document
avec sceau et signature. J’ai l’impression qu’il n’a rien à me dire, et que pour lui, tout
est terminé.

Je lui présente donc mes doléances, selon un ordre progressif, calculé, pour étudier
ses réactions ; Cinquante-six ans après, je résume la substance de cet entretien.

Je signale d’abord la salle de classe qui a besoin d’un bon badigeon, et les coulées de
goudron sur les murs, qui indiquent que les tuyaux de poêle ont besoin d’un bon
ramonage. Il m’écoute, distrait ; ce sont des choses sans importance.

J’en viens alors à la fenêtre de la chambre, branlante et sans vitres. Oh, alors il réagit
sèchement :
-« Monsieur l’instituteur, ne savez vous pas que votre école, c’est un temple qui
appartient à la paroisse de Meyrueis. Adressez vous donc au propriétaire, au Pasteur.
La mairie n’est en rien concernée ! ».
Je réplique aussitôt car la moutarde me monte au nez

« Monsieur le Maire, en tant qu’instituteur public, je me dois de l’ignorer. Je n’ai pas
à en discuter avec le Pasteur, pas plus d’ailleurs qu’avec le Curé. Je m’adresse à
vous, responsable d’une école publique, du logement dû à l’instituteur, de l’entretien
de l’école, je ne suis rien d’autre qu’un instituteur public et laïque, j’ai des droits et
vous, vous avez des devoirs vis à vis de l’école publique. Je n’ai pas à m’immiscer
dans les débats qui vous opposent à la paroisse protestante. Ma neutralité est de
rigueur, mais je ne dois pas être la victime. Vous devriez comprendre cela, Monsieur
le Maire.

J’aborde alors vaillamment, mais sans espoir, la question des wc :
« Peut-on, Monsieur le Maire, concevoir une école sans cour de récréation et sans
WC ? »
J’ai droit alors à une longue péroraison, dont je résume l’essentiel :

« Monsieur l’instituteur, vous avez eu de nombreux prédécesseurs aux Oubrets. Je
n’ai connu que les derniers, mais aucun d’entre eux ne s’est montré aussi exigeant et
aussi ouvertement revendicatif. Ils se sont adaptés à un état de fait, auquel je ne veux
rien changer.
On vous livrera un gros voyage de bois, pour chauffer la salle de classe. ».
J’enregistre la précision : « Scié ? Non ».

Inutile de prolonger la discussion. Je ne l’écoute plus ; je voudrais lui faire sentir mon
indignation et surtout mon mépris. J’aperçois des fioles sur les étagères de la
pharmacie. Alors brusquement :
« donnez moi un tube d’aspirine. »
